\chapter{Introduction}
\label{ch:intro}

Un réseau routier est, par nature, un graphe : les intersections en sont les
sommets et les tronçons de route les arêtes, pondérées par leur longueur.
Calculer le meilleur trajet entre deux points revient alors à chercher un plus
court chemin dans ce graphe, l'un des problèmes les plus classiques et les plus
utiles de l'algorithmique, résolu de manière optimale depuis 1959 par
l'algorithme de Dijkstra.

\medskip

En théorie, le problème est donc « réglé ». En pratique, tout se joue sur
l'\emph{efficacité} : une carte réelle compte des dizaines à des centaines de
milliers d'intersections, et un moteur de routage doit répondre en quelques
millisecondes, souvent pour un très grand nombre de requêtes. L'écart entre
« sait-on le résoudre ? » et « sait-on le résoudre vite, à grande échelle ? » est
précisément l'objet de ce rapport, qui construit et compare, sur des données
routières réelles, toute une échelle d'algorithmes de plus court chemin.

\section{Contexte et motivation}

Trouver le plus court chemin routier entre deux points est sans doute
l'application la plus quotidienne de l'algorithmique : chaque requête adressée à
un GPS ou à un service de cartographie en ligne déclenche, en quelques
millisecondes, la résolution d'un problème de plus court chemin sur un graphe
comptant des millions d'arcs. Derrière la navigation grand public, les mêmes
calculs sous-tendent la logistique du dernier kilomètre, les plateformes de
transport à la demande, la planification de tournées et les services d'urgence, autant de domaines où le temps de réponse et le passage à l'échelle sont
critiques.

\medskip

À cette échelle, l'algorithme de Dijkstra « pur » ne suffit plus : bien qu'il
soit optimal, il explore une part énorme du réseau pour chaque requête, ce qui le
rend trop lent lorsqu'un service doit traiter des milliers de requêtes par seconde
sur un graphe continental. C'est ce constat qui a motivé, depuis les années 2000,
toute une famille de techniques d'\emph{accélération} (heuristiques dirigées
vers le but, recherche bidirectionnelle, pré-calcul par landmarks ou par
raccourcis) qui rendent la \emph{même} réponse optimale en explorant une
fraction infime du graphe. Ce sont ces techniques que le présent projet
implémente, mesure et compare sur des données routières réelles.

\section{Problème et question de recherche}
\label{sec:question}

Tous les algorithmes optimaux considérés rendent, par définition, \emph{la même}
route la plus courte. Ce qui les distingue, c'est le \emph{travail} qu'ils
accomplissent pour y parvenir. La question centrale du projet est donc :

\begin{quote}
\itshape
À réponse optimale identique, dans quelle mesure une meilleure heuristique et de
meilleures structures de données réduisent-elles le nombre de nœuds explorés et
le temps de réponse ?
\end{quote}

La difficulté fondamentale est \textbf{l'échelle}. Une ville réelle compte de
quelques milliers à plusieurs centaines de milliers d'intersections. La stratégie
de base (explorer vers l'extérieur dans toutes les directions jusqu'à
rencontrer le but) gaspille alors un travail considérable : elle visite tout un
disque de nœuds là où la route optimale n'en traverse qu'une fine bande. Toute la
valeur des algorithmes avancés consiste à concentrer l'exploration là où elle est
utile, au prix d'un \emph{pré-calcul} effectué une seule fois puis amorti sur
toutes les requêtes suivantes.

\section{Objectifs et contributions}

Le projet apporte les contributions suivantes :

\begin{bulletlist}
  \item \textbf{Une modélisation rigoureuse} du réseau routier réel en graphe
        orienté pondéré (chapitre \\~\ref{ch:modelisation}), avec les choix
        d'ingénierie (composante fortement connexe, listes d'adjacence
        compactes) qui rendent l'expérimentation possible à l'échelle.
  \item \textbf{Une échelle complète d'algorithmes} (chapitre~\ref{ch:algos}),
        des fondamentaux (BFS, DFS, Dijkstra, Greedy, A*, ALT) aux techniques
        avancées (recherche bidirectionnelle, files à seaux de Dial, landmarks
        actifs, Arc Flags, Contraction Hierarchies), toutes décrites avec leur
        justification mathématique.
  \item \textbf{Une évaluation expérimentale approfondie}
        (chapitre~\ref{ch:xp}) : un banc d'essai agrégé, une étude de l'effet du
        placement des landmarks, une \emph{étude de passage à l'échelle}
        démontrant que les gains croissent avec la taille du graphe, et une
        \emph{analyse statistique} de la difficulté des requêtes.
  \item \textbf{Une mise en perspective} (chapitre~\ref{ch:discussion}) : un
        récapitulatif des notions mises en œuvre, un retour sur l'algorithme IDA*
        essayé puis écarté, et un parallèle entre les landmarks d'ALT et les
        Pattern Databases.
\end{bulletlist}

\section{Correction garantie par construction}
\label{sec:correctness}

Un point méthodologique guide tout le projet : pour que la comparaison ait un
sens, il faut comparer des algorithmes qui rendent \emph{la même} réponse.
Comparer le nombre de nœuds explorés par deux algorithmes n'a en effet d'intérêt
que s'ils calculent tous deux le plus court chemin exact ; sans cette garantie,
un algorithme pourrait « gagner » simplement en rendant une route plus courte à
calculer mais fausse.

\medskip

Chaque algorithme censé être optimal est donc contrôlé, sur \emph{chaque} requête
du banc d'essai, contre la longueur de chemin de Dijkstra pris comme référence :
les requêtes sont tirées d'une graine fixe et sont identiques pour tous les
algorithmes, et le harnais d'expérimentation s'interrompt à la moindre
divergence de longueur. Les tableaux de résultats rapportent d'ailleurs l'erreur
relative maximale observée : elle vaut exactement zéro pour tous les algorithmes
optimaux, ce qui certifie que les écarts de performance mesurés portent bien sur
le \emph{travail} et non sur la qualité de la réponse. Elle n'est non nulle que
pour les algorithmes dont on sait qu'ils sacrifient l'optimalité (Greedy, BFS,
DFS), qui sont alors analysés à part, précisément pour quantifier le prix de cet
abandon.
